% arara: quickstart

% Arara, the cool TeX automation tool
% Copyright (c) 2012 -- 2020, Island of TeX
% All rights reserved.
%
% Redistribution and  use in source  and binary forms, with  or without
% modification, are  permitted provided  that the  following conditions
% are met:
%
% 1. Redistributions  of source  code must  retain the  above copyright
% notice, this list of conditions and the following disclaimer.
%
% 2. Redistributions in binary form  must reproduce the above copyright
% notice, this list  of conditions and the following  disclaimer in the
% documentation and/or other materials provided with the distribution.
%
% 3. Neither  the name  of the  project's author nor  the names  of its
% contributors may be used to  endorse or promote products derived from
% this software without specific prior written permission.
%
% THIS SOFTWARE IS  PROVIDED BY THE COPYRIGHT  HOLDERS AND CONTRIBUTORS
% "AS IS"  AND ANY  EXPRESS OR IMPLIED  WARRANTIES, INCLUDING,  BUT NOT
% LIMITED  TO, THE  IMPLIED WARRANTIES  OF MERCHANTABILITY  AND FITNESS
% FOR  A PARTICULAR  PURPOSE  ARE  DISCLAIMED. IN  NO  EVENT SHALL  THE
% COPYRIGHT HOLDER OR CONTRIBUTORS BE  LIABLE FOR ANY DIRECT, INDIRECT,
% INCIDENTAL, SPECIAL, EXEMPLARY,  OR CONSEQUENTIAL DAMAGES (INCLUDING,
% BUT  NOT LIMITED  TO, PROCUREMENT  OF SUBSTITUTE  GOODS OR  SERVICES;
% LOSS  OF USE,  DATA, OR  PROFITS; OR  BUSINESS INTERRUPTION)  HOWEVER
% CAUSED AND  ON ANY THEORY  OF LIABILITY, WHETHER IN  CONTRACT, STRICT
% LIABILITY, OR TORT (INCLUDING NEGLIGENCE OR OTHERWISE) ARISING IN ANY
% WAY  OUT  OF  THE USE  OF  THIS  SOFTWARE,  EVEN  IF ADVISED  OF  THE
% POSSIBILITY OF SUCH DAMAGE.

\pdfobjcompresslevel=1
\pdfpageattr {/Group << /S /Transparency /I true /CS /DeviceRGB>>}

\documentclass[a4paper,oneside,11pt]{memoir}

\usepackage[T1]{fontenc}
\usepackage[utf8]{inputenc}
\usepackage[margin=2.5cm]{geometry}
\usepackage{arara}

\usepackage{tikzlings}

\newcommand{\araraversion}{%
  \InputIfFileExists{version.tex}{}{undefined}}

\newcommand{\peter}[1][1]{\begin{tikzpicture}[scale=#1, transform shape]
\node[
  draw=araracolour,
  thick,
  fill=araracolour!20!white,
  rectangle,
  minimum width=4cm,
  minimum height=3cm,
  rounded corners
] (box) {};

\begin{scope} 
\clip (box.south west) rectangle (box.north east);
\anteater[
  yshift=-3.2cm,
  xshift=-.3cm,
  scale=2,
  hat=warningcolour
]
\end{scope}
\end{tikzpicture}}

\newcommand{\accessories}{%
\begin{tikzpicture}[
  y=0.80pt,
  x=0.80pt,
  yscale=-.044,
  xscale=.044,
  inner sep=0pt,
  outer sep=0pt
]
% hat
\path[
  fill=black,
  line width=0.080pt
]
(332.9381,0.2962) .. controls (310.4381,2.1962) and (301.5381,3.0962) ..
(289.9381,4.7962) .. controls (211.3381,16.4962) and (164.7381,42.8962) ..
(142.9381,88.0962) .. controls (131.0381,112.7962) and (127.8381,130.5962) ..
(125.9381,181.6962) .. controls (125.4381,197.3962) and (124.4381,217.1962) ..
(123.8381,225.6962) .. controls (122.5381,245.4962) and (119.2381,269.7962) ..
(117.6381,271.3962) .. controls (117.3381,271.6962) and (112.3381,270.9962) ..
(106.5381,269.7962) .. controls (85.2381,265.4962) and (64.2381,266.0962) ..
(46.0381,271.2962) .. controls (30.7381,275.7962) and (22.5381,281.8962) ..
(13.0381,296.0962) .. controls (-1.7619,317.9962) and (-3.9619,338.7962) ..
(6.4381,358.5962) .. controls (15.4381,375.5962) and (45.7381,403.5962) ..
(70.2381,417.5962) .. controls (85.8381,426.3962) and (97.2381,430.6962) ..
(133.4381,441.2962) .. controls (169.9381,451.8962) and (171.7381,452.3962) ..
(192.2381,455.0962) .. controls (276.5381,466.4962) and (311.0381,469.0962) ..
(373.9381,468.9962) .. controls (420.8381,468.9962) and (440.8381,467.8962) ..
(477.4381,463.1962) .. controls (576.5381,450.5962) and (634.4381,433.0962) ..
(672.8381,404.1962) .. controls (694.7381,387.5962) and (708.5381,368.3962) ..
(714.0381,346.6962) .. controls (718.7381,328.2962) and (714.3381,312.9962) ..
(699.0381,295.0962) .. controls (678.4381,270.7962) and (647.1381,260.3962) ..
(612.5381,266.2962) .. controls (607.5381,267.0962) and (600.9381,267.6962) ..
(597.9381,267.4962) -- (592.4381,267.1962) -- (591.8381,262.6962) ..
controls (591.4381,260.1962) and (590.6381,248.6962) .. (590.0381,237.1962) ..
controls (589.4381,225.5962) and (588.2381,203.7962) .. (587.4381,188.6962) ..
controls (584.0381,126.1962) and (577.3381,96.5962) .. (560.7381,71.0962) ..
controls (554.2381,61.0962) and (537.7381,44.5962) .. (527.9381,38.1962) ..
controls (500.3381,20.1962) and (458.5381,7.4962) .. (409.4381,2.1962) ..
controls (396.3381,0.6962) and (343.1381,-0.6038) .. (332.9381,0.2962) -- cycle;

% glasses
\path[
  fill=black,
  line width=0.080pt
] (141.9381,516.6962) .. controls (122.6381,518.5962) and (104.4381,521.6962) ..
(60.9381,530.5962) .. controls (50.2381,532.7962) and (39.0381,534.8962) ..
(36.2381,535.2962) .. controls (33.3381,535.5962) and (30.4381,536.5962) ..
(29.8381,537.2962) .. controls (29.1381,538.1962) and (28.8381,543.1962) ..
(29.1381,552.3962) .. controls (29.5381,569.4962) and (30.7381,572.2962) ..
(41.6381,583.8962) .. controls (51.4381,594.3962) and (53.8381,598.8962) ..
(58.4381,615.5962) .. controls (64.4381,637.4962) and (76.8381,670.5962) ..
(85.3381,687.5962) .. controls (96.3381,709.3962) and (117.6381,722.2962) ..
(155.8381,730.3962) .. controls (167.4381,732.8962) and (170.3381,733.0962) ..
(191.9381,733.0962) .. controls (210.5381,733.0962) and (217.5381,732.7962) ..
(225.4381,731.2962) .. controls (261.7381,724.5962) and (285.6381,708.1962) ..
(308.4381,674.1962) .. controls (317.1381,661.1962) and (325.9381,643.2962) ..
(329.8381,630.4962) .. controls (331.6381,624.7962) and (334.2381,616.4962) ..
(335.5381,612.0962) .. controls (338.6381,602.0962) and (343.2381,594.7962) ..
(347.9381,592.4962) .. controls (353.1381,589.8962) and (371.8381,589.9962) ..
(377.6381,592.5962) .. controls (385.1381,596.0962) and (389.1381,602.7962) ..
(397.9381,627.6962) .. controls (418.0381,684.3962) and (441.0381,710.1962) ..
(484.8381,724.6962) .. controls (517.3381,735.4962) and (548.1381,736.1962) ..
(580.4381,726.8962) .. controls (612.4381,717.5962) and (633.1381,701.2962) ..
(646.4381,674.6962) .. controls (651.3381,664.6962) and (654.5381,655.1962) ..
(662.9381,625.1962) .. controls (669.9381,599.9962) and (675.6381,586.7962) ..
(681.4381,582.0962) .. controls (683.4381,580.4962) and (686.6381,577.7962) ..
(688.5381,575.9962) -- (692.0381,572.7962) -- (691.7381,556.5962) ..
controls (691.4381,540.5962) and (691.4381,540.3962) .. (688.9381,538.5962) ..
controls (687.5381,537.4962) and (682.8381,535.9962) .. (678.4381,535.2962) ..
controls (674.0381,534.4962) and (663.0381,531.9962) .. (653.9381,529.7962) ..
controls (635.4381,525.2962) and (617.4381,521.7962) .. (595.9381,518.5962) ..
controls (583.9381,516.6962) and (577.0381,516.2962) .. (555.4381,516.2962) ..
controls (525.2381,516.1962) and (520.0381,516.8962) .. (462.9381,527.2962) ..
controls (372.7381,543.7962) and (353.8381,544.2962) .. (288.4381,531.1962) ..
controls (248.4381,523.1962) and (229.0381,520.1962) .. (199.9381,517.7962) ..
controls (182.3381,516.2962) and (152.5381,515.6962) .. (141.9381,516.6962) --
cycle(237.4381,532.2962) ..
controls (256.8381,535.0962) and (276.8381,540.7962) .. (291.4381,547.6962) ..
controls (309.0381,555.8962) and (318.3381,566.0962) .. (321.5381,580.3962) ..
controls (323.0381,587.0962) and (323.0381,589.2962) .. (322.0381,597.1962) ..
controls (320.4381,609.6962) and (314.2381,627.2962) .. (305.0381,646.1962) ..
controls (295.2381,666.1962) and (288.7381,675.9962) .. (278.4381,686.1962) ..
controls (261.9381,702.5962) and (241.9381,712.6962) .. (216.9381,717.2962) ..
controls (202.3381,719.9962) and (172.9381,719.7962) .. (158.3381,716.8962) ..
controls (140.0381,713.1962) and (129.7381,709.1962) .. (118.3381,701.0962) ..
controls (106.1381,692.4962) and (98.8381,682.1962) .. (91.5381,663.5962) ..
controls (80.4381,635.1962) and (75.2381,608.3962) .. (76.1381,584.1962) ..
controls (76.5381,573.1962) and (76.9381,571.2962) .. (79.6381,565.4962) ..
controls (85.0381,553.8962) and (92.0381,548.0962) .. (107.0381,542.5962) ..
controls (124.8381,536.0962) and (148.5381,532.8962) .. (197.9381,530.2962) ..
controls (208.1381,529.7962) and (226.4381,530.6962) .. (237.4381,532.2962) --
cycle(564.1381,532.2962) ..
controls (597.8381,535.5962) and (620.2381,541.1962) .. (632.6381,549.2962) ..
controls (642.6381,555.9962) and (647.9381,564.1962) .. (650.0381,576.1962) ..
controls (653.6381,596.8962) and (643.6381,647.3962) .. (630.6381,673.6962) ..
controls (626.2381,682.7962) and (624.3381,685.3962) .. (616.4381,693.2962) ..
controls (603.2381,706.5962) and (590.1381,712.5962) .. (564.6381,717.2962) ..
controls (550.4381,719.9962) and (522.7381,719.7962) .. (508.2381,716.8962) ..
controls (480.4381,711.3962) and (463.6381,702.5962) .. (446.5381,684.6962) ..
controls (428.8381,666.1962) and (411.7381,632.7962) .. (405.4381,604.3962) ..
controls (402.9381,593.0962) and (403.6381,580.0962) .. (407.1381,572.4962) ..
controls (413.8381,557.9962) and (429.7381,547.3962) .. (456.4381,539.6962) ..
controls (488.6381,530.3962) and (521.8381,528.0962) .. (564.1381,532.2962) -- cycle;

% mustache
\path[
  fill=black,
  line width=0.080pt
] (294.9381,961.1962) .. controls (277.8381,964.7962) and (259.4381,974.1962) ..
(239.9381,989.1962) .. controls (236.9381,991.4962) and (225.4381,1000.9962) ..
(214.4381,1010.1962) .. controls (190.2381,1030.3962) and (181.1381,1036.8962) ..
(165.8381,1044.6962) .. controls (149.4381,1053.0962) and (137.7381,1056.8962) ..
(126.2381,1057.4962) .. controls (114.7381,1058.0962) and (109.8381,1056.4962) .. 
(104.0381,1050.1962) .. controls (97.1381,1042.6962) and (94.7381,1029.2962) ..
(99.2381,1022.5962) .. controls (102.3381,1017.9962) and (107.6381,1015.2962) .. 
(117.2381,1013.1962) .. controls (127.3381,1011.0962) and (133.4381,1008.4962) ..
(134.8381,1005.7962) .. controls (135.8381,1004.0962) and (135.5381,1003.4962) .. 
(133.5381,1002.1962) .. controls (121.3381,994.1962) and (102.0381,995.0962) ..
(88.8381,1004.2962) .. controls (78.6381,1011.2962) and (68.7381,1025.7962) ..
(65.4381,1038.2962) .. controls (63.7381,1044.8962) and (63.5381,1056.0962) ..
(64.9381,1062.4962) .. controls (66.6381,1069.8962) and (72.0381,1079.5962) ..
(78.1381,1086.1962) .. controls (111.5381,1122.5962) and (206.0381,1138.0962) ..
(290.6381,1121.1962) .. controls (313.6381,1116.4962) and (326.8381,1111.2962) .. 
(341.4381,1100.6962) .. controls (352.2381,1092.8962) and (357.2381,1090.5962) .. 
(364.1381,1089.9962) .. controls (369.7381,1089.3962) and (370.0381,1089.4962) .. 
(377.1381,1094.3962) .. controls (399.2381,1109.5962) and (430.5381,1119.7962) .. 
(472.4381,1125.3962) .. controls (489.0381,1127.5962) and (541.4381,1127.8962) .. 
(555.9381,1125.8962) .. controls (588.2381,1121.3962) and (613.9381,1113.9962) .. 
(631.6381,1104.0962) .. controls (673.5381,1080.6962) and (680.2381,1040.4962) .. 
(647.4381,1009.8962) .. controls (630.3381,993.9962) and (611.1381,990.9962) .. 
(600.5381,1002.7962) .. controls (598.5381,1004.8962) and (596.9381,1007.0962) .. 
(596.9381,1007.6962) .. controls (596.9381,1009.4962) and (602.5381,1011.4962) .. 
(614.7381,1013.9962) .. controls (628.6381,1016.8962) and (633.0381,1019.2962) .. 
(634.9381,1025.2962) .. controls (637.8381,1034.0962) and (632.5381,1043.8962) .. 
(620.2381,1052.2962) .. controls (611.9381,1057.9962) and (605.1381,1059.4962) .. 
(595.6381,1057.8962) .. controls (580.8381,1055.2962) and (572.4381,1049.8962) .. 
(519.9381,1008.7962) .. controls (482.8381,979.8962) and (463.8381,968.9962) .. 
(439.2381,962.5962) .. controls (432.0381,960.6962) and (427.8381,960.2962) ..
(415.9381,960.2962) .. controls (398.7381,960.2962) and (389.4381,962.1962) ..
(373.0381,969.1962) .. controls (367.3381,971.6962) and (362.1381,973.6962) ..
(361.4381,973.6962) .. controls (360.7381,973.6962) and (355.9381,971.6962) ..
(350.8381,969.2962) .. controls (336.7381,962.5962) and (328.6381,960.5962) ..
(313.9381,960.0962) .. controls (305.6381,959.8962) and (299.2381,960.1962) ..
(294.9381,961.1962) -- cycle;
\end{tikzpicture}}

\newcommand{\dad}[1][1]{\begin{tikzpicture}[scale=#1, transform shape]
\node[
  thick,
  rectangle,
  minimum width=4cm,
  minimum height=3cm,
  rounded corners
] (box) {};

\begin{scope}[
  scale=2,
  yshift=-1.7cm,
  xshift=-.13cm
]
\clip (box.south west) rectangle (box.north east);
\anteater
\node at (0, 1.72) {\accessories};
\end{scope}

\node[
  draw=araracolour,
  rectangle,
  rounded corners=2pt,
  fill=araracolour!20!white,
  inner sep=4pt
] at ([yshift=-.53cm]box.south) {\parbox{3.5cm}{\centering\small\em Anteater facts\\
with Peter's dad}};
\end{tikzpicture}}

\definecolor{pinkglasses}{RGB}{128,0,128}

\newcommand{\hair}{%
\begin{tikzpicture}[
  y=0.80pt,
  x=0.80pt,
  yscale=-1.4,
  xscale=1.4,
  inner sep=0pt,
  outer sep=0pt
]
\path[
  fill=black
] (17.0420,6.0910) .. controls (12.7510,7.3280) and (12.2310,11.1080) ..
(11.8050,13.0730) .. controls (11.3790,15.0400) and (10.6560,17.5560) ..
(11.1750,19.1090) .. controls (10.2290,18.4300) and (10.4470,17.3390) ..
(10.4470,17.3390) .. controls (10.4470,17.3390) and (4.7500,25.7270) ..
(14.6650,29.8970) .. controls (13.9870,30.7220) and (10.7380,29.4140) ..
(10.7380,29.4140) .. controls (10.7380,29.4140) and (10.5200,31.2310) ..
(12.7250,30.6750) .. controls (9.1380,33.7790) and (5.3550,26.0690) ..
(5.3550,26.0690) -- (5.8400,28.2530) .. controls (5.8400,28.2530) and (4.5780,28.5910) ..
(2.6880,17.8760) .. controls (0.7980,7.1580) and (8.0710,1.6300) .. (9.8650,0.9510) ..
controls (11.6590,0.2710) and (14.7140,0.0780) .. (16.3140,0.9510) ..
controls (19.9510,-1.8140) and (23.5730,2.2030) .. (24.4110,3.4720) ..
controls (25.2490,4.7430) and (27.9030,8.5160) .. (28.3880,11.0850) ..
controls (28.8720,13.6530) and (31.2480,31.0860) .. (22.2790,28.6870) ..
controls (25.4300,27.7900) and (25.4300,22.1880) .. (25.4300,22.1880) ..
controls (25.4300,22.1880) and (22.9810,24.7100) .. (22.2300,23.6920) ..
controls (25.4300,22.1880) and (25.5520,15.6690) .. (23.5400,12.0060) ..
controls (21.9730,9.0630) and (18.5870,6.7300) .. (17.0420,6.0910) -- cycle;
\end{tikzpicture}}

\newcommand{\glasses}{%
\begin{tikzpicture}[
  y=0.80pt,
  x=0.80pt,
  yscale=-.2,
  xscale=.2,
  inner sep=0pt,
  outer sep=0pt
]
\begin{scope}[
  shift={(-25.07793,-59.30764)}
]
\path[fill=pinkglasses,nonzero rule,line width=0.028pt] (31.6043,59.6053) ..
controls (28.0060,60.1443) and (25.0779,61.4535) .. (25.0779,62.5231) ..
controls (25.0779,62.7567) and (25.9951,63.3250) .. (27.1240,63.7910) ..
controls (31.3253,65.5249) and (32.2359,66.4389) .. (37.1406,73.8473) ..
controls (38.8222,76.3873) and (41.5881,80.0071) .. (43.3400,81.9612) ..
controls (50.2579,89.6764) and (55.7343,92.5720) .. (66.3529,94.1285) ..
controls (72.3149,95.0023) and (75.6172,94.8746) .. (80.8168,93.5693) ..
controls (88.1546,91.7271) and (92.0344,88.9779) .. (95.4564,83.1959) ..
controls (97.6069,79.5623) and (98.5566,78.9626) .. (102.1599,78.9626) ..
controls (105.7229,78.9626) and (106.7220,79.5760) .. (108.7639,83.0195) ..
controls (110.9603,86.7237) and (112.5710,88.6054) .. (114.8246,90.1008) ..
controls (119.1140,92.9470) and (128.2654,95.1434) .. (133.8746,94.6732) ..
controls (146.2571,93.6346) and (153.4679,90.3231) .. (160.5548,82.4198) ..
controls (163.7718,78.8320) and (166.3481,75.3995) .. (168.7894,71.4484) ..
controls (171.2087,67.5326) and (173.2485,65.5979) .. (176.4902,64.1445) ..
controls (178.0071,63.4647) and (179.2418,62.7210) .. (179.2418,62.4878) ..
controls (179.2418,61.7759) and (177.3015,60.6477) .. (174.9732,60.0057) ..
controls (173.0499,59.4755) and (171.3043,59.3912) .. (163.1904,59.4360) ..
controls (148.3865,59.5174) and (144.2463,60.3013) .. (109.9210,69.5183) ..
controls (105.3349,70.7499) and (101.0751,71.0293) .. (97.7502,70.3170) ..
controls (96.4883,70.0465) and (90.7299,68.5574) .. (84.9443,67.0059) ..
controls (70.5510,63.1447) and (62.6551,61.3417) .. (54.6760,60.0932) ..
controls (49.8429,59.3368) and (35.4310,59.0320) .. (31.6043,59.6053) --
cycle(61.9432,66.4778) .. controls (66.9188,67.3121) and (74.3257,69.2894) ..
(82.5807,71.9864) .. controls (90.2713,74.4992) and (92.8077,76.2134) ..
(92.7975,78.8920) .. controls (92.7852,82.1249) and (89.4599,86.1225) ..
(85.0149,88.2491) .. controls (80.2234,90.5411) and (75.4899,91.2706) ..
(69.5985,90.6250) .. controls (64.3421,90.0490) and (61.0930,89.1024) ..
(56.2635,86.7395) .. controls (50.6896,84.0129) and (46.9510,80.6612) ..
(44.9516,76.5989) .. controls (43.2501,73.1417) and (42.9986,70.5195) ..
(44.1972,68.7320) .. controls (45.7894,66.3578) and (48.9257,65.5231) ..
(55.2404,65.7937) .. controls (57.1807,65.8770) and (60.1970,66.1846) ..
(61.9432,66.4778) -- cycle(157.0168,66.4397) ..
controls (161.0198,67.7657) and (161.9222,71.2014) ..
(159.5833,76.2109) .. controls (156.0940,83.6838) and (145.9396,89.5096) ..
(134.4038,90.6564) .. controls (124.3496,91.6562) and (115.1051,87.8921) ..
(112.1802,81.6084) .. controls (110.7116,78.4531) and (111.6849,76.2441) ..
(115.3538,74.4047) .. controls (118.8759,72.6390) and (133.3102,68.2554) ..
(140.2952,66.8298) .. controls (146.5746,65.5482) and (153.8171,65.3792) ..
(157.0168,66.4397);
\end{scope}
\end{tikzpicture}}

\newcommand{\mum}[1][1]{\begin{tikzpicture}[scale=#1, transform shape]
\node[
  thick,
  rectangle,
  minimum width=4cm,
  minimum height=3cm,
  rounded corners
] (box) {};

\begin{scope}[
  scale=2,
  yshift=-1.7cm,
  xshift=-.13cm
]
\clip (box.south west) rectangle (box.north east);
\anteater
\node at (0, 1.55) {\glasses};
\node at (-0.07, 1.6) {\hair};
\end{scope}

\node[
  draw=warningcolour,
  rectangle,
  rounded corners=2pt,
  fill=warningcolour!20!white,
  inner sep=4pt
] at ([yshift=-.53cm]box.south) {\parbox{3.5cm}{\centering\small\em Smart writing tips\\
with Peter's mum}};
\end{tikzpicture}}

\newtcolorbox{talkbox}{
  sidebyside,
  lower separated=false,
  lefthand width=3.3cm,
  sidebyside gap=10pt,
  colback=white,
  boxrule=0pt,
  frame hidden,
  colframe=white,
  left=0pt,
  right=0pt,
  top=0pt,
  bottom=0pt
}

\newtcolorbox{factbox}{
  sidebyside,
  lower separated=false,
  righthand width=3.3cm,
  sidebyside gap=14pt,
  colback=attentioncolour!20!white,
  boxrule=0pt,
  frame hidden,
  colframe=white,
  left=6pt,
  right=6pt,
  top=4pt,
  bottom=4pt,
  width=.8\textwidth
}

\newtcolorbox{tipbox}{
  sidebyside,
  lower separated=false,
  righthand width=3.3cm,
  sidebyside gap=14pt,
  colback=pinkglasses!20!white,
  boxrule=0pt,
  frame hidden,
  colframe=white,
  left=6pt,
  right=6pt,
  top=4pt,
  bottom=4pt,
  width=.8\textwidth
}

\title{A quickstart guide to \arara}
\author{Peter, the anteater, for the Island of \TeX}

\begin{document}

\maketitle

\parskip=1em

Welcome to \arara, the cool \TeX\ automation tool! This document is intended to be quickstart guide, providing the basic instructions to start using our tool. We will help our good friend Peter\footnote{Peter was graduated from Termite High School, Alta Floresta, and is now teaching as a fellow at Mato Grosso University. He is a corresponding member of Duckpond Academy, Sempione Park, Milano, Italy. His doctoral thesis ``On flexibility: \LaTeX, latex and rubber'' was published at Manaus. Some of his scientific works arose from a fruitful cooperation with the well known Brazilian-Italian savant Professore P.\ van Duck.}, the anteater, in his learning adventure. He is very excited about it! Hello, Peter!

\begin{talkbox}
\peter[.7]
\tcblower
\em Hello, nice to meet you, person from the Island of \TeX! I am Peter, the anteater. So I heard about this tool named \arara, what is it? Is there anything to do with \TeX\ and friends?
\end{talkbox}

As a matter of fact, it is! \arara\ is a \TeX\ automation tool. The name was chosen as an homage to a Brazilian bird of the same name, which is a macaw. The word \emph{arara} comes from the Tupian word \emph{a'rara}, which means \emph{big bird}. The tool is an effort to provide a concise way to automate the daily \TeX\ workflow for users and also package writers.

\begin{talkbox}
\peter[.7]
\tcblower
\em It definitely sounds intriguing! So tell me, how does \arara\ work? Does it do everything automatically for me, like other tools? Do things happen by magic?
\end{talkbox}

A very good question, Peter. The best way to explain how \arara\ works is to provide a quick comparison with similar tools. I believe you are familiar with \TeX\ documents, right? Let us use the following file \rbox{article.tex} as an example:

\begin{ncodebox}{Source file}{teal}{\icnote}{white}{article.tex}
\documentclass{article}

\begin{document}
Hello world!
\end{document}
\end{ncodebox}

How would one successfully compile \rbox{article.tex} with \rbox{latexmk} and \rbox{rubber}, for instance? It is quite straightforward: it is just a matter of providing the file to the tool and letting it do the hard work:

\begin{codebox}{Terminal}{teal}{\icnote}{white}
$ latexmk -pdf article.tex
$ rubber --pdf article.tex
\end{codebox}

The mentioned tools perform an analysis on the file and decide what has to be done. However, if one tries to invoke \rbox{arara} on \rbox{article.tex}, I am afraid \emph{nothing} will be generated; the truth is, \arara\ does not know what to do with your file, and the tool will even raise an error message complaining about this issue:

\begin{codebox}{Terminal}{teal}{\icnote}{white}
$ arara article.tex
  __ _ _ __ __ _ _ __ __ _
 / _` | '__/ _` | '__/ _` |
| (_| | | | (_| | | | (_| |
 \__,_|_|  \__,_|_|  \__,_|

Processing "article.tex" (size: 70 B, last modified: 12/28/2020
07:03:16), please wait.

  ERROR

It looks like no directives were found in the provided file. Make
sure to include at least one directive and try again.

Total: 0.04 seconds
\end{codebox}

Peter, this is the major difference of \arara\ when compared to other tools: \emph{it is not an automatic process and the tool does not employ any guesswork on its own}. You are in control of your documents; \arara\ will not do anything unless you \emph{teach it how to do a task and explicitly tell it to execute the task}. So, it is not magic.

\begin{talkbox}
\peter[.7]
\tcblower
\em I see. On the other hand, this approach gives me, as user, full control of my document compilation. Being aware of the compilation steps makes me understand the \TeX\ tooling better. Quite clever! Now, do tell me: how does one tell \arara\ to do a task?
\end{talkbox}

First of all, Peter, \arara\ has to know \emph{how} to do a task. This is done by defining rules. A rule is a formal description of how the tool should handle a certain task. For example, if we want to use \rbox{pdflatex} with \arara, we need a rule for that. Once a rule is defined, the tool automatically provides an access layer to the user. Note that \arara\ is distributed with dozens of predefined rules, so you already have several options out of the box to set up your workflow.

Now, back to your question. Once we know how to execute a task, we need to explicitly tell \arara\ when to do it. That is actually the easy part, provided that you have everything up and running. We accomplish the task by adding a special comment line, hereafter known as \emph{directive}, somewhere in our document to indicate how the tool should work, e.g: 

\begin{codebox}{Source file}{teal}{\icnote}{white}
% arara: pdflatex
\end{codebox}

\begin{talkbox}
\peter[.7]
\tcblower
\em Makes sense. So, a directive is a special comment line which is not the command to be executed, but the name of the rule associated with that directive. Is that correct?
\end{talkbox}

Perfect, Peter! That is basically how \arara\ works: we teach the tool to do a task by providing a rule, and tell it to execute it
via directives in the source code. Let us see how our \rbox{article.tex} file should be:

\begin{ncodebox}{Source file}{teal}{\icnote}{white}{article.tex}
% arara: pdflatex
\documentclass{article}

\begin{document}
Hello world!
\end{document}
\end{ncodebox}

Note, Peter, that the tool expects \emph{you} to provide a list of tasks, and this is done by inserting special comments, i.e, directives, in the source file. Since we just added one directive in our document, let us see how \arara\ behaves with this updated code:

\begin{codebox}{Terminal}{teal}{\icnote}{white}
$ arara article.tex
  __ _ _ __ __ _ _ __ __ _
 / _` | '__/ _` | '__/ _` |
| (_| | | | (_| | | | (_| |
 \__,_|_|  \__,_|_|  \__,_|

Processing "article.tex" (size: 88 B, last modified: 12/28/2020
07:05:05), please wait.

(PDFLaTeX) PDFLaTeX engine .............................. SUCCESS

Total: 0.56 seconds
\end{codebox}

\begin{talkbox}
\peter[.7]
\tcblower
\em Goodness me, it worked like a charm! The moment we specified a directive, \arara\ knew exactly what I wanted to do with that particular file! Awesome!
\end{talkbox}

Exactly! You see, Peter, there is no guesswork from \arara. The tool will do exactly what you tell it to do, no more, no less. There is a lot of freedom to this design, which gives you an interesting way to enhance your \TeX\ experience.

\begin{talkbox}
\peter[.7]
\tcblower
\em I have been wondering: there are scenarios in which we need to provide additional information to the underlying commands~---~for instance, we need to enable shell escape when using the \rbox[araracolour]{minted} package. How can we achieve this? 
\end{talkbox}

For such scenarios, \arara\ provides a second type of directive, a parametrized one, which allows passing arguments to the corresponding rule. From your example, to enable shell escape, one simply needs to write the following directive:

\begin{codebox}{Source file}{teal}{\icnote}{white}
% arara: pdflatex: { shell: yes }
\end{codebox}

Of course, \rbox[cyan]{shell} is not taken randomnly, but defined in the rule scope, otherwise \arara\ would raise an error about an invalid key. The reference manual has a list of all available keys for each predefined rule. It is worth a read.

\begin{talkbox}
\peter[.7]
\tcblower
\em Great, these directives are really convenient! I am now curious on how to explore \arara: shall we move to a more complex document? Consider the following addition to my document:  
\end{talkbox}

\begin{ncodebox}{Source file}{teal}{\icnote}{white}{article.tex}
\documentclass{article}

\begin{document}
\section{Termite recipes}
\label{sec:termiterecipes}

Hello, this is Section~\ref{sec:termiterecipes}.
\end{document}
\end{ncodebox}

\begin{talkbox}
\peter[.7]
\tcblower
\em As we can see, this document has to be compiled twice, or the reference to the first section will not be resolved accordingly. How can I tackle this scenario with \arara?  
\end{talkbox}

The solution is quite straightforward, Peter: how about adding two directives into your document? You can keep them together for convenience (usually at the top), but they can happen anywhere in your file. Let us update the code:

\begin{ncodebox}{Source file}{teal}{\icnote}{white}{article.tex}
% arara: pdflatex
% arara: pdflatex
\documentclass{article}

\begin{document}
\section{Termite recipes}
\label{sec:termiterecipes}

Hello, this is Section~\ref{sec:termiterecipes}.
\end{document}
\end{ncodebox}

The execution workflow is now as expected: our document was correctly compiled twice, as it should be, so references are resolved accordingly! 

\begin{codebox}{Terminal}{teal}{\icnote}{white}
$ arara article.tex 
  __ _ _ __ __ _ _ __ __ _
 / _` | '__/ _` | '__/ _` |
| (_| | | | (_| | | | (_| |
 \__,_|_|  \__,_|_|  \__,_|

Processing 'article.tex' (size: 196 B, last modified: 01/07/2021
08:18:00), please wait.

(PDFLaTeX) PDFLaTeX engine .............................. SUCCESS
(PDFLaTeX) PDFLaTeX engine .............................. SUCCESS

Total: 1.60 seconds
\end{codebox}

\begin{talkbox}
\peter[.7]
\tcblower
\em Cool! But I have been wondering: once the references are resolved, subsequents runs will add an extra, unnecessary compilation. It would not do any harm, surely, but is there a way to avoid it?   
\end{talkbox}

There is a way, Peter! We can use logical expressions and special operators and methods processed at runtime in order to determine whether and how a directive should be processed. This feature is named \emph{directive conditional}, or simply \emph{conditional} for short. Let us update the document to include this feature:

\begin{ncodebox}{Source file}{teal}{\icnote}{white}{article.tex}
% arara: pdflatex
% arara: pdflatex if found('log', 'undefined references')
\documentclass{article}

\begin{document}
\section{Termite recipes}
\label{sec:termiterecipes}

Hello, this is Section~\ref{sec:termiterecipes}.
\end{document}
\end{ncodebox}

Observe, Peter, that the first directive has no conditional, so it will always be processed. The second one, however, has an associated logical test: check if the log file contains a warning about undefined references and, if so, process the directive itself. When references are resolved accordingly, the tool will not process the second directive at all, as we can see in the following run:

\begin{codebox}{Terminal}{teal}{\icnote}{white}
$ arara article.tex 
  __ _ _ __ __ _ _ __ __ _
 / _` | '__/ _` | '__/ _` |
| (_| | | | (_| | | | (_| |
 \__,_|_|  \__,_|_|  \__,_|

Processing 'article.tex' (size: 236 B, last modified: 01/07/2021
08:42:02), please wait.

(PDFLaTeX) PDFLaTeX engine .............................. SUCCESS

Total: 0.97 seconds
\end{codebox}

In this particular case, the test is evaluated beforehand, and the directive is processed if, and only if, the result of such evaluation is true. This directive, when the conditional holds true, is executed at most once.

\begin{talkbox}
\peter[.7]
\tcblower
\em These conditionals can help me dictate how my workflow should behave on certain scenarios! I gather there are several possible ways of tackling them, right? 
\end{talkbox}

Correct, Peter. In the previous example, we used the \rbox[attentioncolour]{\color{black}if} operator, but we could have used \rbox[attentioncolour]{\color{black}while} or \rbox[attentioncolour]{\color{black}until} as well. The logical expression could also be rewritten. We could even have combined these two directives into one! So you see, there is always room for improvement. If this subject caught your attention, take a look at the reference manual for more details on directive conditionals and available methods.

\begin{talkbox}
\peter[.7]
\tcblower
\em I will certainly check them out! Now, I am curious to see other directives in action together, so let us try a different yet common scenario: bibliographies and citations. Consider the following bibliography file, containing a reference to my doctoral thesis: 
\end{talkbox}

\begin{ncodebox}{Source file}{teal}{\icnote}{white}{bibliography.bib}
@phdthesis{peter:2020,
  author  = {Peter Anteater}, 
  title   = {On flexibility: \LaTeX, latex and rubber},
  school  = {Polytechnic University of P{\^{a}}ntano Fundo},
  year    = 2020,
  address = {Pantanal},
  month   = {jan}
}
\end{ncodebox}

\begin{talkbox}
\peter[.7]
\tcblower
\em I want to cite my thesis in the article you and I are writing for this quickstart guide. From the \TeX\ side, it is quite straightforward! How about \arara? Here is my code:
\end{talkbox}

\begin{ncodebox}{Source file}{teal}{\icnote}{white}{article.tex}
\documentclass{article}

\begin{document}
\section{Termite recipes}
\label{sec:termiterecipes}

A citation from my thesis~\cite{peter:2020}.

\bibliographystyle{alpha}
\bibliography{bibliography}
\end{document}
\end{ncodebox}

Peter, remember that \arara\ does not employ any guesswork, so do not expect it to magically compile your document and get all the citations correctly. We need to explicitly tell it what to do, as if we were explaining the compilation steps to a friend. \arara\ is focused on reproducibility. We, as users, can learn a great deal about \TeX\ and friends by organising our workflow into directives. Let us review what should be done in order to correctly compile your document:

\begin{enumerate}
\item For starters, we need to run the \TeX\ engine in order to write (amongst other things) the relevant bibliography information to the auxiliary file. We can achieve this by inserting the following directive:

\begin{codebox}{Directive}{teal}{\icnote}{white}
% arara: pdflatex
\end{codebox}

\item Once the auxiliary file holds the relevant bibliography information, we need to run the Bib\TeX\ tool as a means to map entries from the bibliography database to the existing citations into an intermediate bibliography file. We can achieve this by inserting the following directive:

\begin{codebox}{Directive}{teal}{\icnote}{white}
% arara: bibtex
\end{codebox}

\item Once the intermediate bibliography file is generated, we can now have a proper bibliography section in our document, so we need to run the \TeX\ engine again. We can achieve this by inserting the following directive:

\begin{codebox}{Directive}{teal}{\icnote}{white}
% arara: pdflatex
\end{codebox}

\item However, the citations are still not yet referenced in the document, so a final run of the \TeX\ engine is required in order to correctly adjust these references. We can achieve this by inserting the following directive:

\begin{codebox}{Directive}{teal}{\icnote}{white}
% arara: pdflatex
\end{codebox}
\end{enumerate}

And we are done! These are the compilation steps required to correctly generate your article from the given source and bibliography files. Note that the entire workflow involves two different tools working together: the \TeX\ engine and the Bib\TeX\ tool.

\begin{talkbox}
\peter[.7]
\tcblower
\em Wait a minute, no less than four compilation steps? I would never guessed it! This is great: \arara\ is also helping me understand better how the \TeX\ workflow works!
\end{talkbox}

Quite true, Peter! \arara\ tries its best to help users think out of the box! Now that we know the exact compilation steps to be taken, we just need to arrange them in the correct order in our document:

\begin{ncodebox}{Source file}{teal}{\icnote}{white}{article.tex}
% arara: pdflatex
% arara: bibtex
% arara: pdflatex
% arara: pdflatex
\documentclass{article}

\begin{document}
\section{Termite recipes}
\label{sec:termiterecipes}

A citation from my thesis~\cite{peter:2020}.

\bibliographystyle{alpha}
\bibliography{bibliography}
\end{document}
\end{ncodebox}

\end{document}
